\documentclass[11pt]{article}
\usepackage{docstyle}
\renewcommand\sheetname{Light \textperiodcentered{} In-class work answers}
\FigAnstrue

\begin{document}
\sheethead{Light: In-class Work \textperiodcentered{} Answers}{Tutor's working, not an official mark scheme}{Answers only}
Drawings show the given parts in black and the answer in \ans{red}.

\section*{Part A \quad Shadows and colour}
\A{A1} Two straight rays from the lamp that just touch the top and bottom of the ball, continued to the screen, with arrows. The region behind the ball between them is the \textbf{umbra}. A small lamp gives an umbra only: no penumbra.
\begin{center}\TaskShadow[5.6mm]\end{center}
\A{A2} white paper in green light: \textbf{green} \quad blue car in white light: \textbf{blue} \quad blue car in red light: \textbf{black} \quad yellow banana in green light: \textbf{green} (yellow objects reflect red and green) \quad black shoe in white light: \textbf{black}
\A{A3} red + blue = \textbf{magenta}; \quad green + blue = \textbf{cyan}

\section*{Part B \quad Reflection}
\begin{minipage}[t]{0.5\textwidth}
\vspace{0pt}
\A{B1} $i = 90^\circ - 30^\circ = \mathbf{60^\circ}$; $r = \mathbf{60^\circ}$ (law of reflection). The normal is drawn at 90\textdegree{} to the mirror; both angles are marked from the normal.
\end{minipage}\hfill
\begin{minipage}[t]{0.46\textwidth}
\vspace{0pt}\centering\TaskReflect[5.8mm]{0}
\end{minipage}

\begin{minipage}[t]{0.5\textwidth}
\vspace{0pt}
\A{B2} Each image point is on a line at 90\textdegree{} to the mirror, the \textbf{same distance behind} it as the object point is in front. A$'$B$'$C$'$ is drawn dotted (virtual image).
The two rays: draw straight lines from A$'$ to the top and bottom of the eye; where they cross the mirror are the reflection points. Join A to those points (real rays, arrows), and show the parts behind the mirror dotted.
\end{minipage}\hfill
\begin{minipage}[t]{0.46\textwidth}
\vspace{0pt}\centering\TaskTriangle[6mm]
\end{minipage}

\begin{minipage}[t]{0.58\textwidth}
\vspace{0pt}
\A{B3} Two \textbf{parallel} mirrors at \textbf{45\textdegree}. The ray from the object hits the top mirror, turns through 90\textdegree{} down the tube, hits the bottom mirror and turns through 90\textdegree{} into the eye. At each mirror $i = r = 45^\circ$.
\end{minipage}\hfill
\begin{minipage}[t]{0.38\textwidth}
\vspace{0pt}\centering\FigPeriscope[5.4mm]{1}
\end{minipage}

\section*{Part C \quad Refraction}
\begin{minipage}[t]{0.5\textwidth}
\vspace{0pt}
\A{C1} Normals (dashed) at both faces. Inside the block the ray bends \textbf{towards} the normal (angle of incidence 35\textdegree, angle in the glass about 22\textdegree). At the bottom face it bends \textbf{away} from the normal and leaves \textbf{parallel} to the incident ray, shifted sideways.
\end{minipage}\hfill
\begin{minipage}[t]{0.46\textwidth}
\vspace{0pt}\centering\FigBlock[5.8mm]{1}
\end{minipage}

\begin{minipage}[t]{0.5\textwidth}
\vspace{0pt}
\A{C2} The ray goes \textbf{from the fish} to the surface and bends \textbf{away from the normal} into the eye (arrows pointing towards the eye). Extend the ray in air straight back as a dotted line: the fish appears \textbf{higher} (closer to the surface), above and slightly towards the viewer. The real fish is \textbf{below} where it appears.
\end{minipage}\hfill
\begin{minipage}[t]{0.46\textwidth}
\vspace{0pt}\centering\TaskFish[5.6mm]{Sp}
\end{minipage}

\A{C3} \textbf{P} is red. Each colour travels at a different speed in glass, so each is refracted by a different amount. Red is refracted the \textbf{least}, so it is the ray bent least from its original direction (P); violet is bent most (Q).

\section*{Part D \quad The eye and lenses}
\A{D1} A \textbf{cornea} \quad B \textbf{iris} \quad C \textbf{pupil} \quad D \textbf{lens} \quad E \textbf{retina} \quad F \textbf{optic nerve} \quad G \textbf{blind spot}

\begin{minipage}[t]{0.5\textwidth}
\vspace{0pt}
\A{D2} The three rays \textbf{converge} and meet at the focal point F on the axis; the middle ray does not bend. The focal length is the distance from the centre of the lens to F.
\end{minipage}\hfill
\begin{minipage}[t]{0.46\textwidth}
\vspace{0pt}\centering\FigConvex[5.6mm]{2}
\end{minipage}

\A{D3} The blind spot is where the optic nerve leaves the eye. There are \textbf{no light-sensitive cells} there, so light that lands on it is not detected and no signal is sent to the brain.

\end{document}
